\begin{ex}
Ước lượng tích $49 \times 21$ bằng cách làm tròn các thừa số đến hàng chục: $50 \times 20 = \dots$
\shortans{1000}
\loigiai{
$50 \times 20 = 1000$.
}
\end{ex}

\begin{ex}
Ước lượng tích $38 \times 52$ bằng cách làm tròn đến hàng chục: $40 \times 50 = \dots$
\shortans{2000}
\loigiai{
$40 \times 50 = 2000$.
}
\end{ex}

\begin{ex}
Ước lượng tích $69 \times 31$ bằng cách làm tròn đến hàng chục: $70 \times 30 = \dots$
\shortans{2100}
\loigiai{
$70 \times 30 = 2100$.
}
\end{ex}

\begin{ex}
Ước lượng tích $78 \times 42$: làm tròn thành $80 \times 40 = \dots$
\shortans{3200}
\loigiai{
$80 \times 40 = 3200$.
}
\end{ex}

\begin{ex}
Ước lượng thương $812 : 39$ bằng cách làm tròn: $800 : 40 = \dots$
\shortans{20}
\loigiai{
$800 : 40 = 20$.
}
\end{ex}

\begin{ex}
Ước lượng thương $594 : 28$ bằng cách làm tròn: $600 : 30 = \dots$
\shortans{20}
\loigiai{
$600 : 30 = 20$.
}
\end{ex}

\begin{ex}
Ước lượng thương $896 : 31$ bằng cách làm tròn: $900 : 30 = \dots$
\shortans{30}
\loigiai{
$900 : 30 = 30$.
}
\end{ex}

\begin{ex}
Một cửa hàng bán 21 chiếc xe đạp, mỗi chiếc giá khoảng 1 950 000 đồng (khoảng 2 triệu đồng). Ước lượng tổng số tiền thu được là bao nhiêu triệu đồng?
\shortans{42}
\loigiai{
$21 \times 2 = 42$ triệu đồng.
}
\end{ex}

\begin{ex}
Một rạp chiếu phim có 29 hàng ghế, mỗi hàng có 31 ghế. Ước lượng rạp có khoảng bao nhiêu ghế bằng cách làm tròn đến hàng chục?
\shortans{900}
\loigiai{
$30 \times 30 = 900$ ghế.
}
\end{ex}

\begin{ex}
Mỗi ngày bạn Nam đọc khoảng 19 trang sách. Ước lượng trong 31 ngày Nam đọc được bao nhiêu trang sách?
\shortans{600}
\loigiai{
$20 \times 30 = 600$ trang.
}
\end{ex}

\begin{ex}
Một thùng nước chứa khoảng 198 lít nước. Ước lượng 5 thùng như thế chứa khoảng bao nhiêu lít nước?
\shortans{1000}
\loigiai{
$200 \times 5 = 1000$ lít.
}
\end{ex}

\begin{ex}
Một người đi bộ mỗi phút được khoảng 78 m. Ước lượng trong 60 phút người đó đi được bao nhiêu mét?
\shortans{4800}
\loigiai{
$80 \times 60 = 4800$ m.
}
\end{ex}

\begin{ex}
Làm tròn số 4872 đến hàng trăm được số nào?
\shortans{4900}
\loigiai{
Làm tròn lên 4900.
}
\end{ex}

\begin{ex}
Làm tròn số 5129 đến hàng nghìn được số nào?
\shortans{5000}
\loigiai{
Làm tròn thành 5000.
}
\end{ex}

\begin{ex}
Ước lượng tổng $4980 + 3020$ bằng cách làm tròn đến hàng nghìn: $5000 + 3000 = \dots$
\shortans{8000}
\loigiai{
$5000 + 3000 = 8000$.
}
\end{ex}

\begin{ex}
Ước lượng hiệu $7950 - 2040$ bằng cách làm tròn đến hàng nghìn: $8000 - 2000 = \dots$
\shortans{6000}
\loigiai{
$8000 - 2000 = 6000$.
}
\end{ex}

\begin{ex}
Một quyển sách có 192 trang, đã đọc 89 trang. Ước lượng còn khoảng bao nhiêu trang chưa đọc?
\shortans{100}
\loigiai{
$190 - 90 = 100$ trang.
}
\end{ex}

\begin{ex}
Một xe buýt chở được 48 người. Ước lượng 10 xe buýt chở được bao nhiêu người?
\shortans{500}
\loigiai{
$50 \times 10 = 500$ người.
}
\end{ex}

\begin{ex}
Ước lượng tích $102 \times 48$: lấy $100 \times 50 = \dots$
\shortans{5000}
\loigiai{
$100 \times 50 = 5000$.
}
\end{ex}

\begin{ex}
Ước lượng thương $418 : 21$: lấy $400 : 20 = \dots$
\shortans{20}
\loigiai{
$400 : 20 = 20$.
}
\end{ex}

\begin{ex}
Một mảnh đất hình chữ nhật dài 39 m, rộng 19 m. Ước lượng diện tích mảnh đất khoảng bao nhiêu mét vuông?
\shortans{800}
\loigiai{
$40 \times 20 = 800\text{ m}^2$.
}
\end{ex}

\begin{ex}
Một trang trại có 98 chuồng lợn, mỗi chuồng nuôi 11 con. Ước lượng trang trại nuôi khoảng bao nhiêu con lợn?
\shortans{1000}
\loigiai{
$100 \times 10 = 1000$ con.
}
\end{ex}

\begin{ex}
Tính chính xác tích $50 \times 20 = \dots$
\shortans{1000}
\loigiai{
$50 \times 20 = 1000$.
}
\end{ex}

\begin{ex}
Làm tròn 98 đến hàng chục được số nào?
\shortans{100}
\loigiai{
Làm tròn thành 100.
}
\end{ex}

\begin{ex}
Làm tròn 103 đến hàng trăm được số nào?
\shortans{100}
\loigiai{
Làm tròn thành 100.
}
\end{ex}
